\section{Experimental Setup}
\label{sec:setup}

\paragraph{Data and splits.}
We use HybridQA \citep{chen2020hybridqa}, whose questions need a table and linked Wikipedia passages.
The official test labels are never opened.
We build five splits from the official train and dev files (Table~\ref{tab:splits}): \emph{train} (3{,}000 questions, a seeded subsample stratified by answer-node type) and \emph{dev} (1{,}000 questions; used for all tuning and model selection) from official train tables, with disjoint tables; \emph{test} from official-dev questions on tables absent from train and dev; and two out-of-domain (OOD) splits from \emph{held-out table clusters}.
Tables are clustered with $k$-means ($k=20$, seed 0, fitted on train-table embeddings of title, section title and header with the same encoder); clusters 6, 16 and 19 are held out entirely.
\emph{ood\_test} uses official-dev questions on held-out-cluster tables, and \emph{ood\_dev} (500 questions) is kept for diagnostics and never used for tuning.
A leakage audit asserts disjoint table ids across all pairs of splits and disjoint clusters for the OOD splits.
Test and ood\_test are evaluated once per final configuration.
HybridQA has no annotated hop order, so gold paths are heuristic: we take the weak \texttt{answer-node} label, add the same-row anchor cell with the highest question overlap, and append the linked passage when the answer is in a passage.
Questions without an answer node have no gold path ($1.7\%$ overall).
No path is verified by hand; we report this noise in Section~\ref{sec:limitations} and Appendix~\ref{app:data}.

\begin{table}[t]
\centering\small
\begin{tabular}{@{}lrrr@{}}
\toprule
Split & Questions & Tables & Gold path (heur.) \\
\midrule
train & 3{,}000 & 2{,}564 & 2{,}972 \\
dev & 1{,}000 & 844 & 988 \\
test & 2{,}045 & 1{,}787 & 1{,}991 \\
ood\_test & 569 & 498 & 547 \\
ood\_dev & 500 & 429 & 495 \\
\bottomrule
\end{tabular}
\caption{HybridQA splits used here. The last column counts questions with a heuristic gold path; the rest have none.}
\label{tab:splits}
\end{table}

\paragraph{Methods.}
All path-based methods use the same encoder, candidate generator, policy architecture (where a policy is trained) and frozen answer head.
\textbf{M0} (oracle) feeds the answer head the gold-path evidence and bounds what path-based methods can reach.
\textbf{B1} (SFT) fine-tunes \texttt{flan-t5-small} on the question and a BM25-ranked linearization of table rows with linked-passage snippets, without hops.
\textbf{B2a} is trace-level GRPO-style RL: the same policy and action space, $G=8$ sampled trajectories per question, trajectory reward equal to answer F1, advantages normalized within each group and applied to every action, a clipped objective as in PPO and no value head.
\textbf{B2b} is the supervision control: B2a with the trace-level reward $\mathrm{F1}(\hat a,a^\ast)+0.2\cdot\mathrm{F1}_{\text{set}}(\tau,g)$, which uses the gold path but only after the episode ends.
\textbf{B3} is a differentiable soft selector: an MLP scores hop-0 candidates and their one-hop neighbours, is trained by marginal likelihood of candidates whose text contains the gold answer (no path labels), and at inference passes its top 3 candidates to the answer head.
\textbf{H1} is \hop\ with the full hybrid reward.
Ablations \textbf{A1}--\textbf{A4} are defined in Section~\ref{sec:method}.

\paragraph{Matched budgets.}
The budget unit is the number of (question, trajectory) rollouts processed during training: 48{,}000 for every RL method (GRPO counts $G$ per question: 8 questions $\times$ 8 samples per update), and 48{,}000 question passes (16 epochs over the 3{,}000 training questions) for B1 and B3.
Tuning is also matched (Appendix~\ref{app:changes}): each trained method receives 3 trials on dev with seed 0 over a fixed grid (Appendix~\ref{app:hp}), and reward weights are not tuned; ablations inherit the setting chosen for H1.
We report wall-clock cost in Appendix~\ref{app:hp}.
The trials and the selected setting per method are listed in Table~\ref{tab:study} (Appendix~\ref{app:hp}); one RL run takes 15--40 minutes and one B1 run about 75 minutes on our laptop.

\paragraph{Metrics.}
\emph{EM} and token-level \emph{F1} use the HybridQA normalization (lowercase, no punctuation or articles), with the maximum over gold answers.
Path metrics use only questions with a gold path.
\emph{Strict path validity} is the fraction of selected hops that are legal, not redundant and on the gold path; \emph{loose validity} drops the gold-path condition.
\emph{Gold-path agreement} is the set F1 between selected entities and the gold path, and an ordered variant (longest common subsequence divided by gold length).
\emph{Hop efficiency} is $\min(1,|g|/n_{\text{hops}})$ if the answer is exactly correct, and $0$ otherwise.
\emph{Mean hops} is the number of selected entities.
B1 has no hops, so path metrics are not defined for it.

\paragraph{Seeds and statistics.}
Each final configuration is trained with 3 seeds.
We report the mean and standard deviation over seeds, and a 95\% percentile bootstrap interval (10{,}000 resamples over questions) on pooled per-question correctness.
Comparisons of H1 with each other method use a paired bootstrap over the same questions (difference in EM, 95\% interval, two-sided $p$).
With three seeds we do not rely on seed-level tests, and we describe a difference as supported only when the criteria of the claims ledger are met.
The ID$\rightarrow$OOD gap is $\mathrm{EM}_{\text{test}}-\mathrm{EM}_{\text{ood\_test}}$ (absolute) and that divided by $\mathrm{EM}_{\text{test}}$ (relative).
The seeds column of Tables~\ref{tab:main} and~\ref{tab:ablations} gives the number of completed seeds per method (\R{SeedsHOneHoppo} for \hop\ in the current build); no completed run was dropped. Runs trained before the design changes of Appendix~\ref{app:changes} are excluded and archived.

\paragraph{Compute and implementation.}
All runs use compute tier S: one Apple M4 laptop (10 cores, 16\,GB unified memory), CPU for the policy and MPS for the frozen encoder and answer head, and no GPU.
Backbones are small and frozen (bge-small encoder, flan-t5-small answer head); only a $\approx$0.3M-parameter policy (or, for B1, the flan-t5-small model) is trained.
We make no claim about larger models.
Configurations, seeds and the code version of every run are logged, and all tables and figures are generated from the run logs (Appendix~\ref{app:repro}).

\paragraph{Synthetic sanity study.}
A separate, clearly labelled study (E6) runs H1 on programmatically generated tables with exact gold paths and a rule-based answer head; it checks that the learning signal can be exploited without label noise and is never presented as a benchmark result.
On synthetic tables \hop\ reaches \SanityEm\% EM and \SanityValid\% strict path validity, so the hop-level signal is learnable when gold paths are exact.
